\documentclass[lang=cn,newtx,10pt,sCheme=Chinese,citestyle=gb7714-2015, bibstyle=gb7714-2015]{../../Template/elegantbook}

% citestyle=gb7714-2015, bibstyle=gb7714-2015是为了将参考文献格式设置为国标GB7714-2015,不想使用这个文献格式标准的话可以删去

%%%%%%%%%%%%%%%%%%%%%%%%%%%%%%%%%导言区%%%%%%%%%%%%%%%%%%%%%%%%%%%%%%%%%%%%%%%%%%

\title{数学分析}

\author{实空}
\institute{无}
\date{\today}
\version{ElegantBook-4.5}
\bioinfo{全局显示/隐藏证明/解/笔记}{\AllToggle}% 全局显示/隐藏证明/解答/笔记

\extrainfo{宠辱不惊,闲看庭前花开花落;
\\
去留无意,漫随天外云卷云舒.}

\setcounter{tocdepth}{3}

\logo{blank-photo.png}
\cover{cover.png}

% 修改标题页的橙色带
\definecolor{customcolor}{RGB}{32,178,170}
\colorlet{coverlinecolor}{customcolor}

% 本文档额外使用的宏包和命令

\usepackage{../../Styles/mystyle-elegantbook}% 注意根据文件位置修改模板和宏包的索引！！！

% 交叉引用设置（这里不要引用这个文件自身）
% 注意这里的引用地址不是.sty文件相对于其他被引用的.tex文件地址,而是需要引用的.tex文件相对于其他被引用的.tex文件地址!!!
% 条件加载：仅当目标 .aux 文件存在时才引入，避免因缺少文件而报错
\IfFileExists{../Abstract Algebra/main.aux}{\externaldocument[Abstract Algebra-]{../Abstract Algebra/main}}{}
\IfFileExists{../Analytic Geometry/main.aux}{\externaldocument[Analytic Geometry-]{../Analytic Geometry/main}}{}
\IfFileExists{../Basis of Algebra/main.aux}{\externaldocument[Basis of Algebra-]{../Basis of Algebra/main}}{}
% \IfFileExists{../Basis of Analytics/main.aux}{\externaldocument[Basis of Analytics-]{../Basis of Analytics/main}}{}
\IfFileExists{../Complex Analysis/main.aux}{\externaldocument[Complex Analysis-]{../Complex Analysis/main}}{}
\IfFileExists{../Differential Geometry/main.aux}{\externaldocument[Differential Geometry-]{../Differential Geometry/main}}{}
\IfFileExists{../Functional Analysis/main.aux}{\externaldocument[Functional Analysis-]{../Functional Analysis/main}}{}
\IfFileExists{../Harmonic Analysis/main.aux}{\externaldocument[Harmonic Analysis/-]{../Harmonic Analysis//main}}{}
\IfFileExists{../Numerical Analysis/main.aux}{\externaldocument[Numerical Analysis-]{../Numerical Analysis/main}}{}
\IfFileExists{../Probability Theory/main.aux}{\externaldocument[Probability Theory-]{../Probability Theory/main}}{}
\IfFileExists{../Real Analysis/main.aux}{\externaldocument[Real Analysis-]{../Real Analysis/main}}{}
\IfFileExists{../Set Theory/main.aux}{\externaldocument[Set Theory-]{../Set Theory/main}}{}
\IfFileExists{../Topology/main.aux}{\externaldocument[Topology-]{../Topology/main}}{}

% 注意如果上述跨文档交叉引用代码报错，先执行Tasks: Run Task命令删除所有辅助文件，再编译

\usepackage{subfiles}% 主文件额外引用宏包 尽量置于导言区的最后


%%%%%%%%%%%%%%%%%%%%%%%%%%%%%%%%%导言区%%%%%%%%%%%%%%%%%%%%%%%%%%%%%%%%%%%%%%%%%%

\begin{document}

%%%%%%%%%%%%%%%%%%%%%%%%%%%%正文格式设置%%%%%%%%%%%%%%%%%%%%%%%%%%%%%%%%%%%%%%%%%%

\maketitle
\frontmatter

\tableofcontents% 添加目录

\mainmatter% 将行为改回预期版本，并重置页码

%%%%%%%%%%%%%%%%%%%%%%%%%%%%正文格式设置%%%%%%%%%%%%%%%%%%%%%%%%%%%%%%%%%%%%%%%%%%


%%%%%%%%%%%%%%%%%%%%%%%%%%%%%%%%%%正文部分%%%%%%%%%%%%%%%%%%%%%%%%%%%%%%%%%%%%%%%%

% \subfile{contents/chapter-1.tex}

\subfile{contents/chapter-3.tex}

\subfile{contents/chapter-4.tex}

\subfile{contents/chapter-11.tex}

\subfile{contents/chapter-5.tex}

\subfile{contents/chapter-6.tex}

\subfile{contents/chapter-15.tex}

\subfile{contents/chapter-13.tex}

\subfile{contents/chapter-14.tex}

\subfile{contents/chapter-16.tex}

\subfile{contents/chapter-17.tex}

\subfile{contents/chapter-7.tex}

\subfile{contents/chapter-8.tex}

\subfile{contents/chapter-12.tex}

\subfile{contents/chapter-2.tex}

% \subfile{contents/chapter-10.tex}

\subfile{contents/chapter-19.tex}

\appendix

\subfile{contents/chapter-18.tex}

\subfile{contents/chapter-9.tex}

% \subfile{contents/chapter-20.tex}

% \subfile{contents/chapter-21.tex}

% \subfile{contents/chapter-22.tex}


%%%%%%%%%%%%%%%%%%%%%%%%%%%%%%%%%%正文部分%%%%%%%%%%%%%%%%%%%%%%%%%%%%%%%%%%%%%%%%

%%%%%%%%%%%%%%%%%%%%%%%%%%%%%%%%%%参考文献%%%%%%%%%%%%%%%%%%%%%%%%%%%%%%%%%%%%%%%%

\nocite{*}% 显示所有参考文献（原本只显示已引用的参考文献）
\printbibliography[heading=bibintoc, title=\ebibname]

%%%%%%%%%%%%%%%%%%%%%%%%%%%%%%%%%%参考文献%%%%%%%%%%%%%%%%%%%%%%%%%%%%%%%%%%%%%%%%

\end{document}

